\label{sec:3bound4}

To extend the results of the previous section to the case of 3--manifolds bounding 4--manifolds, we consider generic proper maps from closed orientable 3--manifolds into the 2--ball $B^2$.
Such a map is Morse--like in the sense that its singular set is well behaved, it can be studied via the same techniques as in Section \ref{sec:2bound3}, and we may recover a Stein factorization and Stein complex that define a handle decomposition for a bounding 4--manifold.

\begin{theorem}
	\label{thm:3bound4}
	Let $M$ be a closed orientable 3--manifold and $f:M\to \BB$ a proper generic map.
	Then there exists a cobordism of the pair $(M,\emptyset)$ with an explicit handle decomposition described fully by an integer decorated Stein complex of $f$.
\end{theorem}

\begin{proof}%[Proof of Theorem \ref{thm:3bound4}]
	By Sard's theorem, the image of the singular set $f(s_f)$ in $\BB$ has Lebesgue measure 0.
	By genericity, and because the maps constructed in Chapter \ref{cha:alg1} satisfy this requirement, we take $s_f$ to consist of a set of arcs in the plane that intersect each other only pairwise and transversally, and in such a way that $\BB\setminus s_f$ consists of a collection of regions $R_i$ homeomorphic to copies of $\BB$ along with a single annular region $R_\infty$ that does not intersect $f(M)$.
	This leaves us with a classification of critical values of $f$ into codimension 1 critical values, i.e.\ arcs away from crossings, and codimension 2 critical values, i.e.\ arc crossings.
	This classification gives $f(s_f)$ a natural structure as a 4--valent planar graph, where the vertex set is given by the set of codimension 2 critical values and the edge set is the collection of codimension 1 critical values.
	

	We obtain a 4--manifold with boundary $M$ by attaching 2--handles corresponding to the regular regions $R_i$, 3--handles corresponding to codimension 1 singularities, and 4--handles corresponding to codimension 2 singularities.
	
	We will first focus on a region $R$ in $\BB$ that is not $R_\infty$.
	Begin by ``shrinking'' $R$ away from $f(s_f)$ into a space $\D$ that is homeomorphic to $\DD$
	Every point $x\in D$ has preimage a disjoint collection of copies of $S^1$ by the regular value theorem, giving the preimage $f\inv(\D)$ the structure of a disjoint collection of 1--sphere bundles over $\DD$.
	As $\DD$ is contractible, these bundles are trivial, and the connected components of $f\inv(\D)$ are solid tori.
	We indexed the regions by $R_i$, so the disc $\D$ is indexed $\D_i$ and the $k\nth$ solid torus that maps over $D_i$ is denoted $V_{i,k}$.
	
	\begin{figure}[h]
		\caption{A closed sleeve around an arc of codimension 1 critical values}
		\centering
		\includegraphics[height=3cm]{figures/dummy.jpg}
		\label{fig:arcsleeve}
	\end{figure}
	
	Next, consider an arc $\alpha$ of codimension 1 critical values that separates a pair of regions $R_i$ and $R_j$ where either $i$ or $j$ may be $\infty$.
	As when we considered regions, $\alpha$ has been shrunk away from the codimension 2 critical values it connects.
	Let $\varepsilon(\alpha)$ be a closed ``sleeve'' around $\alpha$ whose boundary intersects the boundary of each of $D_i$ and $D_j$ as in Figure \ref{fig:arcsleeve}.
	The restriction of $f$ to the preimage of a linear transversal to $\alpha$ that connects $D_i$ and $D_j$ is, itself, a Morse function.
	The intersection with $\alpha$ is a critical point of this restriction, and the same techniques from the proof of Theorem \ref{thm:2bound3} may be applied to $f\inv(\varepsilon(\alpha))$.
	We find that $f\inv(\varepsilon(\alpha))$ is the union of either a single copy of $P\times\I$, the interval crossed with a pair of pants, or a single copy of $E\times\I$, the interval crossed with a cup or cap extremum, with a disjoint collection of $A\times\I$, the interval crossed with an annulus.
	In particular, whenever one of the regions that $\alpha$ separates is $R_\infty$, we get the situation where the singular fiber over a point in $\alpha$ is a single point corresponding to a cup/cap extremum.
	When $\varepsilon(\alpha)$ pulls back to a pair of pants crossed with an interval, $\alpha$ is called an \emph{indefinite fold}.
	Otherwise, $\varepsilon(\alpha)$ pulls back to a cup/cap crossed with an interval and $\alpha$ is a \emph{definite fold} in this case.
	
	\begin{figure}[h]
		\caption{A closed sleeve around a codimension 2 critical value}
		\centering
		\includegraphics[height=2cm]{figures/dummy.jpg}
		\label{fig:isolatesleeve}
	\end{figure}
	
	Let $x$ be a codimension 2 critical value $x$ and $\varepsilon(x)$ its neighbourhood that intersects the boundaries of the shrunken regions and arc sleeves it is near as in Figure \ref{fig:isolatesleeve}.
	The possible singular fibers over $x$ are cataloged in \cite{Saeki}, and Figure \ref{fig:saekising} displays them.
	When the singular fibers are disconnected, $f\inv(\varepsilon(x))$ is categorized by the codimension 1 analysis.
	Otherwise, we can find a full analysis in Section 4.4 of \cite{CostThur08}.
	In each case, $\varepsilon(x)$ pulls back to a genus 3 handlebody whose boundary is $f\inv(\pd\varepsilon(x))$.
	It is shown in \cite{CostThur08} that filling in $f\inv(\pd\varepsilon(x))$ with 2--discs and gluing the two genus 3 handlebodies together results in a Heegard splitting of $S^3$.
		
	\begin{figure}
		\caption{Possible singular fibers of a proper generic smooth map from an orientable 3--manifold to a surface}
		\centering
		\includegraphics[height=4cm]{figures/dummy.jpg}
		\label{fig:saekising}
	\end{figure}
	
	A full analysis of this situation is considered in , and the critical preimage to such a sleeve has multiple forms.
	The singular fibers over $x$ are catalogued 
	
	
	











		
\end{proof}



Return now to the case of a closed, oriented 3--manifold $M$ and the generic smooth map $f:M\to \RR$.
It is assumed that $M$ is compact, so we will assume that the image $f(M)$ is also compact.
To build a 4--manifold $W$ whose boundary is exactly $M$, we take the 4--manifold with boundary $M\times I$.
We again let $M\times \{0\}=M_-$ be the boundary component with orientation opposite that of $M$, and $M\times\{1\}=M_+$ the boundary component with orientation exactly that of $M$.
For reasons that will become clear later, we fill in the boundary component $M_+$, setting the domain of $f$ to be $M_+$.
This will result in a 4--manifold whose boundary is $M_-$.
The image of the singular set of $f$ is, as stated before, a set of arcs in the plane.
These arcs separate the plane into connected regions of regular values, and we can choose $f$ so that these regions are all homeomorphic to discs.
The discs of regular values all pull back through $f$ to open solid tori in $M$.
We fill in $M_+$ by attaching 2--, 3--, and 4--handles where the attachment instructions are obtained through $f$.

	First, we take care of the regular values of $f$.
	Let $R_0$ be a disc of regular values, and shrink the region away from the critical values of $f$ slightly.
	This ``shrunken'' region is topologically closed, and homeomorphic to the closed 2--dimensional disc.
	Name this ``shrunken region'' by $R$, and take $f\inv (R)$ to be a collection of attaching neighbourhoods for 4--dimensional 2--handles.
	Take a pair of arbitrary points $p$, $q$ interior to $R$, and let $T$ be an open solid torus that projects over $R$ through $f$.
	Let the circle $f\inv (p)\cap T$ be an attaching sphere for a 2--handle.
	The orientation of $\RR$ and the curve $f\inv (q)\cap T$ determines a framing of $f\inv (p)$, and this pair completely determines the 2--handle attached over $f\inv (R)$.

Next, we extend over the codimension 1 singularities.
An arc $a_0$ of codimension 1 critical values are the image of an arc $s_0$ of critical points of $f$ in $M_+$, and the connected component of $f\inv (a_0)$ containing $s_0$ is an interval crossed with the singularities that occur in the 2--dimensional case.
The endpoints of the arc $a_0$ are codimension 2 critical values of $f$, and the endpoints of the arc $s_0$ are critical points that map to the endpoints of $a_0$.
Slightly shrink $a_0$ away from its boundary to the arc $a$ of critical values of $f$.
Take $A$ to be a tubular neighbourhood of $a$ in $\RR$.
Then $A$ is a copy of $I\times I$ so that lines in the $x$ direction of $A$ are parallel to $a$ and lines in the $y$ direction are normal to $a$.
The boundary components $I\times \{0\}$ and $I\times \{1\}$ run parallel to $a$ and sit in the boundary of the shrunken regions to either side of $a_0$, and the normal boundary components $\{0\}\times I$ and $\{1\}\times I$ normally intersect $a_0$ at $\pd a$.
For a given normal arc $\{x\}\times I$, our analysis of the singularities of the Morse function one dimension down allow us to say that $f\inv(\{x\}\times I)$ is a disjoint collection of cylinders with either a pair of pants or a disc.
Each point of $(I\times \{0\})\cup (I\times\{1\})$ pulls back to a circle that is a boundary component of one of these three shapes.
Pulling $I\times \{0\}$ and $I\times\{1\}$ back through $f$ then yields cylinders.
The 2--handles attached over the preimages of the shrunken regions to either side of $a_0$ have filled these cylinders with discs.
Boundary circles of the $f\inv (\{x\}\times I)$ are filled by discs in $(M\times I)\cup (2\text{--handles})$, so $f$ pulls the normal arcs of $A$ back to spheres in $(M\times I)\cup (2\text{--handles})$.
This means that each connected component of $f\inv (A)$ is a copy of $S^2\times D^1$, to which we attach 4--dimensional 3--handles.

Finally, we extend over codimension 2 singularities.
The codimension 1 singularities can be classified into definite or indefinite folds.
This classification depends on the type of 1--dimensional Morse singularity found by pulling an arc in the plane, normal to the given singular arc, back through $f$ and examining $f$ restricted to the connected component of this preimage containing the critical point that maps through $f$ to the given singular arc.
If the component we've found is a pair of pants, then the singular arc is an indefinite fold.
If it is a disc, then the arc is a definite fold.
Every codimension 2 critical value is either the crossing of a pair of definite folds, a pair of indefinite folds, or of one definite and one indefinite fold.
A codimension 2 critical is mapped to by exactly two critical points in $M$.
Put $x$ to be our codimension 2 critical point, and $p_1$, $p_2$ to be the critical points of $f$ projecting over $x$.
If $x$ is the crossing of a pair of definite folds then we can analyze $p_1$, $p_2$ using Morse theory, so $p_1$, $p_2$ are each of index either 0 or 2.
We get that the connected components of $f\inv (x)$ containing $p_1$, $p_2$ are just the points $p_1$, $p_2$.
If $x$ is the crossing of one definite and one indefinite fold, then, without loss of generality, $p_1$ is of index 1 and $p_2$ is index 0 or 2.
Then the connected component containing $p_1$ is the figure--eight graph from above and the component containing $p_2$ is just $p_2$.
When $x$ is the crossing of two indefinite folds, each of $p_1$, $p_2$ is of index 1.
Our analysis only guarantees that $p_1$, $p_2$ are the centres of crosses with arms alternating in and out of $p_i$ in the connected components of $f\inv(x)$ containing $p_i$.
If $p_{1,2}$ are in different connected components, then each of those components is a figure--eight graph.
If $p_{1,2}$ are in the same connected component, there are exactly two different directed graphs that contain exactly two vertices, each of degree four, that each have two edges entering and two leaving.

Each crossing is in the boundary of exactly four regions of regular values and in the boundary of exactly four singular arcs.
Label the regions $R_1,\dots,R_4$ and the singular arcs by $a_{1,2},a_{2,3}, a_{3,4}, a_{4,1}$.
The regions are labeled in an anticlockwise order about $x$, and the arcs are named for the pair of regions they border.
Note that every point in $\RR$ outside of the range of $f$ is vacuously a regular value of $f$.
Let $x$, $p_{1,2}$ be as they were above.
Let $D_x$ be a small disc in the plane containing $x$.
There is an obvious triangulation of $D_x$ with exactly four triangles so that one vertex of each triangle corresponds to $x$, the shared boundary of a pair of triangles corresponds to an arc $a_{i,i+1}$, and the interior of each triangle corresponds to a region $R_i$ of regular values whose boundary contains $x$.
Let $R_i$ be any of the regions containing $x$ so that $f\inv (R_i)$ is nonempty.
At this point in our construction, we may form an arc $\gamma_i$ inside of $R_i$ from $a_{i-1,i}$ to $a_{i,i+1}$ that consists entirely of points in the boundary of the shrunken region pulled back through $f$ to attach 2--handles and points in the boundary of the tubular neighbourhood of the shrunken arc pulled back through $f$ to attach 3--handles.
If every $R_i$ is in the image of $f$, then the union of the $\gamma_i$ is the boundary of $D_x$.
If some region $R_i$ is outside of the image of $f$, then the arcs $a_{i-1,i}$ and $a_{i,i+1}$ are each definite folds, and the regions to either side of $R_i$ are in the image of $f$.
This is a special case where $x$ is the crossing of a pair of definite folds, and $x$ sits on the boundary of $f(M)\subset\RR$, and in this case we do not define a $\gamma_i$.

Each $D_x$ pulls back through $f$ to a collection of disjoint 3--handlebodies that, together with the boundaries of handles already attached in $(M\times I)\cup$({2-- and 3--handles}), form copies of $S^3$ to which we attach 4--dimensional 4--handles.
When all 4--handles are attached, the component $M_+$ is completely filled, and we have a 4--manifold whose boundary is exactly $M$.
We prove this first for the connected components that consist entirely of regular points, then use similar arguments to prove this for all other cases.

Let $U$ be a component of $f\inv (D_x)$ that consists entirely of regular points.
Then the connected component of $f\inv (x)$ that sits inside of $U$ is a circle, and $T$ is a closed regular neighbourhood of this circle, i.e.\ a solid torus in $M$.
A meridian of this solid torus maps through $f$ to the boundary of $D_x$.
Through handle attachment we have filled every circle in $M_+$ that maps to a point in $\pd D_x$ with a disc.
Then $U$ shares a boundary with a solid torus $V$ in $(M\times I)\cup$(2-- and 3--handles), and the meridian of $U$ is a longitude of $V$, so $U\cup V\subset (M\times I)\cup$(2-- and 3--handles) is a 3--sphere.
We use this 3--sphere as an attaching neighbourhood of a 4--handle.

Next, we consider the case where $U$ contains only one of $p_1$ or $p_2$, abbreviated to $p$, which we recall as being a critical point of $f$ so that $f(p)=x$.
If $p$ is a definite fold, then a pair of opposite singular arcs, say $a_{4,1}$ and $a_{2,3}$, are also definite folds, and the image of $f$ near the strand of critical points $f\inv (a_{4,1}\cup x \cup a_{2,3})$ lies entirely on one side of $a_{4,1}\cup x \cup a_{2,3}$ in the plane, which we will take to be $R_1\cup a_{1,2}\cup R_2$ without loss of generality.
The strand $a_{1,2}$ pulls back through $f$ to only regular points in $U$, so these are circles.
It is easy to see that $U$ is a 3--ball, and $f\inv (\gamma_1\cup\gamma_2)$ is a 2--sphere filled with discs to make another 3--ball.
Gluing these 3--balls together results in another 3--sphere over which we may attach a 4--handle.

If $p$ is an indefinite fold, then the component of $f\inv (x)$ containing $p$ is an oriented figure--eight graph embedded in $U$, $U$ is a regular neighbourhood of the figure--eight graph, and $f$ maps $\pd U$ to the boundary of $D_x$.
Again, the boundary of $U$ maps through $f$ to 

In each remaining case $U$ shares a boundary with another handlebody of the same genus and the two form a Heegard Diagram of $S^3$.
The proof is covered in detail in Section 4.4 of \cite{CostThur08}.
At this point, $M\cup$(2--, 3--, and 4--handles) has only one boundary component: $M_-$.
We obtain a 4--manifold whose boundary is $M$, and whose handle decomposition contains handles of index at most 2, by building the handle decomposition dual to that we've just described.
The process can be simplified somewhat.
For example, the only 4--handles that can't be dealt with by extending 2-- and 3-- handles are those that were added when $p_1$ and $p_2$ were in the same connected component of $f\inv (x)$.
The object of study that simplifies our work considerably is called the \emph{Stein complex} of the Morse 2--function $f:M\to\RR$.

Let $f$ be a Morse function or Morse 2--function with compact fibers as before.
Define the \emph{Stein factorization} of $f$ to be a factorization $f=g\comp h$ so that the following are satisfied:
\begin{enumerate}
  \item The map $h$ has connected fibers, and
  \item the map $g$ is finite--to--one.
\end{enumerate}
We can also see the image of $h$ as the quotient space of the domain of $f$ by the connected fibers of $f$.
The image of $h$ will also be called the \emph{Stein complex} of the map $f$.

In the case of the oriented surface $\Sigma$, the Stein complex of $f$ is a graph whose vertices correspond to the singularities of $f$.
Vertices corresponding to critical points of index 0 or 2 being of degree 1, and vertices corresponding to critical points of index 1 being of degree 3.
We form a 3--manifold whose boundary is $\Sigma$ by starting with a disc bundle over the edges of $h(\Sigma)$, then extending over the vertices by attaching 3--balls as in the analysis above.
There is an obvious deformation retraction of the resulting 3--manifold to the Stein complex of $f$.

In the case of the 3--manifold $M$, the Stein complex is generically a surface and the possible singularities are more interesting.
As in the analysis above, some singularities are just the singularities from the lower dimensional Morse function crossed with an interval, so the corresponding shape in the Stein complex is simply a degree--one or degree--three vertex crossed with an interval.
The more interesting parts of the Stein complex are brought forth at the crossing of indefinite folds in the plane.
As in the case of $\Sigma$, $M$ is a circle bundle over the Stein complex at generic points and the 4--manifold $W$ with $\pd W = M$ is a disc bundle.
We extend over singularities by attaching handles and, as before, $W$ retracts onto the Stein complex.
The Stein complex, along with some extra information we can obtain from $f$, provides a set of building instructions for $W$ explored in the next section.

\begin{cor}
	\label{cor:3stein4}
	Let $M$ be a closed, orientable 3--manifold.
	Then there exists a Morse 2--function $f:M\to\RR$ with Stein factorization $f=g\comp h$ and a Stein complex $S=h(M)$ whose regions can be assigned framing constants such that $S$ provides a full set of instructions for the construction of a 4--manifold $W$ with $\pd W=M$.
\end{cor}







For $M$ a closed, oriented 3--manifold, we would like to describe a 4--manifold $W$ with $\pd W=M$.
An appropriate Morse 2--function $f:M\to\DD$ with Stein complex $S$ is a starting point.
The preimage circles of the regular values of $f$ can all be filled in to copies of $\DD$.
Extending over the 

In the case of the oriented surface $\Sigma$, the Stein complex of $f$ is a graph whose vertices correspond to the singularities of $f$.
Vertices corresponding to critical points of index 0 or 2 being of degree 1, and vertices corresponding to critical points of index 1 being of degree 3.
We form a 3--manifold whose boundary is $\Sigma$ by starting with a disc bundle over the edges of $h(\Sigma)$, then extending over the vertices by attaching 3--balls as in the analysis above.
There is an obvious deformation retraction of the resulting 3--manifold to the Stein complex of $f$.

In the case of the 3--manifold $M$, the Stein complex is generically a surface and the possible singularities are more interesting.
As in the analysis above, some singularities are just the singularities from the lower dimensional Morse function crossed with an interval, so the corresponding shape in the Stein complex is simply a degree--one or degree--three vertex crossed with an interval.
The more interesting parts of the Stein complex are brought forth at the crossing of indefinite folds in the plane.
As in the case of $\Sigma$, $M$ is a circle bundle over the Stein complex at generic points and the 4--manifold $W$ with $\pd W = M$ is a disc bundle.
We extend over singularities by attaching handles and, as before, $W$ retracts onto the Stein complex.
The Stein complex, along with some extra information we can obtain from $f$, provides a set of building instructions for $W$ explored in the next section.


\begin{cor}
	Let $M$ be a closed, orientable 3--manifold.
	Then there exists a Morse 2--function $f:M\to\RR$ with Stein factorization $f=g\comp h$ and a Stein complex $S=h(M)$ whose regions can be assigned framing constants such that $S$ provides a full set of instructions for the construction of a 4--manifold $W$ with $\pd W=M$.
\end{cor}






From the regular value theorem, the preimage of a regular value of $f$ is the disjoint union of oriented circles in $M$.
The image of the singular set of $f$ is a collection of arcs in $\RR$ that only intersect pairwise and transversally.
This leaves us with a classification of critical values of $f$ into codimension 1 singularities, i.e.\ arcs away from crossings, and codimension 2 singularities, i.e.\ arc crossings.
